\problemname{Bitflip}

This is a run-twice problem.

In the first pass, you will be given $n$ bit-strings $S_1, S_2, \ldots, S_n$.
For each string $S_i$, you are allowed to modify it by prepending or appending a \texttt{0} or \texttt{1}.
You may also leave the string unchanged.
Let $T_i$ denote the string modified by your program from $S_i$.

Then, for each string $T_i$, the judge may flip all bits in it (\texttt{0} becomes \texttt{1} and vice versa) or leave the string unchanged.
Let $T_i'$ denote the string modified by the judge.

In the second pass, you will receive the strings $T_i$.
You must recover each original string $S_i$ correctly.

\subsection*{Interaction}

Your program will be run twice. The time and memory limits are applied independently for each of the two runs.

\subsubsection*{First Run}

The first line of the input contains a string \texttt{first} and an integer $n$ ($1 \leq n \leq 10\,000$).
Next, $n$ lines follow.
The $i$-th line contains a bit-string $S_i$ ($1 \leq |S_i| \leq 20$).

Your program must output $n$ lines.
The $i$-th of them contains the bit-string $T_i$.

\subsubsection*{Second Run}

The first line of the input contains a string \texttt{second} and an integer $n$.
Next, $n$ lines follow, each line containing one bit-string $T_i'$.

Your program must output $n$ lines.
The $i$-th of them contains the bit-string $S_i$.
